\documentclass[12pt]{article}
\usepackage[T1]{fontenc}
\usepackage[utf8]{inputenc}
\usepackage{lmodern}
\usepackage{textcomp}
\usepackage{enumitem}
\usepackage{geometry}
\usepackage{graphicx}
\usepackage{amsmath, amssymb}
\geometry{margin=1in}
\begin{document}
\begin{center}
Sociology - Graduate Level Assessment 25 \
Total Marks: 40 \
Answer all questions.
\end{center}

\section*{Multiple Choice (10 marks)}
\begin{enumerate}[label=\arabic*.]
\item 'Scientific' theories in the nineteenth century tried to explain race in naturalistic terms. Which of the following ideas was not considered?
\begin{enumerate}[label=(\alph*)]
    \item genetics
    \item evolution
    \item height
    \item brain size
\end{enumerate}
\item Which of these was not one of Townsend's (1979) indicators of relative deprivation?
\begin{enumerate}[label=(\alph*)]
    \item lacking exclusive use of a bath or shower
    \item living in housing with serious structural defects
    \item buying fewer than twenty DVDs in the previous year
    \item going without a week's holiday in the previous year
\end{enumerate}
\item Weber (1919) said that the state's monopoly of the use of force was legitimated by:
\begin{enumerate}[label=(\alph*)]
    \item charismatic authority
    \item rational-legal authority
    \item traditional authority
    \item value-rational authority
\end{enumerate}
\item The market model of state welfare is based on the principle of:
\begin{enumerate}[label=(\alph*)]
    \item individuals buying welfare privately, with some means-tested benefits
    \item regular benefit payments to men as earners of the 'family wage'
    \item a universalist system of welfare for all, regardless of income
    \item decommodifying social welfare through state provision
\end{enumerate}
\item Judith Butler (1999) suggested that:
\begin{enumerate}[label=(\alph*)]
    \item sexual characteristics are the biological determinants of gender
    \item heterosexuality and homosexuality are essential, opposing identities
    \item the 'two-sex' model replaced the 'one-sex' model in the eighteenth century
    \item gender is performed through bodily gestures and styles to create 'sex'
\end{enumerate}
\end{enumerate}

\section*{True / False (10 marks)}
\textit{Answer True or False}
\begin{enumerate}[label=\arabic*.]
\setcounter{enumi}{5}
\item True or False: Thatcher's government actively expanded financial support for single parents, students, and the unemployed during her tenure.
\item True or False: Chodorow (1978) proposed that boys are primarily attached to their fathers while girls are attached to their mothers during socialization.
\item True or False: Researchers often have unrestricted access to official statistics on domestic violence, facilitating accurate assessments.
\item True or False: Decarceration emphasizes community alternatives to imprisonment and institutional care as effective responses to crime.
\item True or False: The human relations approach advocates for strict control and discipline as the primary means to achieve high productivity.
\end{enumerate}

\section*{Long Form Response (20 marks)}
\textit{Answer each question with a clear and complete explanation.}
\begin{enumerate}[label=\arabic*.]
\setcounter{enumi}{10}
\item Discuss the concept of new forms of community within sociology, analyzing the role and characteristics of various communities, and explain why sociological communities formed by unpopular lecturers do not fit this categorization.
\item Discuss Émile Durkheim's concept of social facts and analyze why he believed they should be the primary focus of sociology.
\end{enumerate}

\vfill
\noindent\textit{End of Paper}
\end{document}